\section{Experimental setup}
\label{sec:setup}

\subsection{Corpus}
Our primary corpus is the public \pcorpusName{} trajectory release \cite{yoonholee2026tb2}: 52{,}104 agent
trials over 89 terminal tasks, each with the full step-by-step trace of messages, tool calls and
observations. From it we draw a stratified analysis sample of \pnTraj{} trajectories (\pnTasks{} tasks,
\pcorpusScaffolds{} scaffolds, \pcorpusModels{} models, \pcorpusSteps{} agent steps, overall pass rate
\pcorpusPass{}\%), balanced across scaffold, outcome and trajectory length. A second corpus (SWE-agent
trajectories over SWE-bench \cite{nebius2024sweagenttraj}) was prepared for a cross-corpus check;
Appendix~\ref{app:transfer} explains why it could not be evaluated here. Terminal-Bench tasks are short and
self-contained --- a run averages \pcorpusStepsMean{} steps in our sample (median \pcorpusStepsMedian{},
maximum \pcorpusStepsMax{}) --- which is both an advantage and a limitation: it makes annotation tractable
and lets us see complete runs, but the conclusions are strongest for runs of tens of steps.

\subsection{Data quality controls}
The release is a scrape, and two artefacts matter. Many observation fields are redacted to opaque
placeholders: of the \pcorpusRedactSteps{} steps that carry any observation, \pcorpusRedactFrac{}\% show a
bare reference instead of output, and averaged per trajectory the share is \pcorpusRedactMean{}\% --- for
some scaffolds the majority are. And a minority of runs are degenerate, consisting mostly of one repeated
agent message with no tool call at all, a harness-level failure rather than agent behaviour. We exclude
trajectories whose redaction rate exceeds 50\% or whose most frequent message covers more than 40\% of
steps, and report sensitivity to that filter. Redaction hurts the \evid{} and \ver{} channels directly,
since a redacted observation exposes no entities and no state, so it biases every channel against itself
rather than in favour of the proposed one; it also drives annotation disagreement.

\subsection{Annotation}
Labels come from a written codebook (Appendix~\ref{app:codebook}) applied by readers who see the whole
trajectory; monitors see only the prefix. Three rounds were collected:

\begin{itemize}[leftmargin=1.2em,itemsep=2pt,topsep=2pt]
\item \textbf{Round 1 (prevalence).} \pannPositionWindows{} windows drawn at evenly spaced positions from
40 trajectories stratified by scaffold, outcome and length. Because the positions are not chosen by any
detector, the positive rate of this sample estimates the prevalence of stagnation
(Section~\ref{sec:results}).
\item \textbf{Round 2 (stability).} Every window of round 1 re-labelled by a \emph{different} reader who
saw only the window and three steps of context on either side, giving a conservative estimate of label
stability under reduced information.
\item \textbf{Round 3 (density).} \pannDenseWindows{} windows from 51 trajectories, roughly ten per
trajectory at most four steps apart, so that consecutive windows of the same label merge into contiguous
stagnant and productive \emph{regions}. Alarm-level metrics are computed on these regions.
\end{itemize}

Each window is labelled by one reader per round; a window read in two rounds is adjudicated by majority,
and any conflict becomes \textsc{uncertain} and is held out of the binary task. That task takes
\textsc{stagnant} and \textsc{done\_redundant} as positive, \textsc{productive} and \textsc{regression} as
negative, and excludes and counts \textsc{blocked\_external} and \textsc{uncertain}. In total
\pannWindowsTotal{} windows over \pannTrajTotal{} trajectories are labelled, of which \pannBinaryN{} enter
the binary task (\pannPosRate{}\% positive).

\subsection{Protocol}
Features are computed for every step of every sampled trajectory and for two window sizes ($w=\pwParam$
and $w=20$). Evaluation uses \pnFolds-fold task-level cross-validation: in each fold, monitors are
trained and standardised on the training tasks and scored on the held-out tasks, so every trajectory is
evaluated by a model that never saw its task. All reported window metrics pool the folds; all alarm
metrics pool the folds as well, and we report trajectory-level bootstrap intervals.

The alarm policy is fixed: a monitor alarms at the first step $t$ where its score is at least $\theta$
for two consecutive steps and $t$ exceeds a two-step margin. We sweep $\theta$ over $[0.30,1.00]$ in
steps of $0.02$ and report, for each false-stop budget, the threshold with the best net savings.
Thresholds are selected on the annotated set, which is the correct protocol for a deployment decision but
does mean the reported operating point is optimistic; the full curve is therefore reported alongside it.

\subsection{Metrics}
\emph{Window level}: ROC-AUC, PR-AUC (average precision) with the positive prevalence marked, best-F1 with
its threshold, and the mean ROC-AUC computed \emph{inside} each trajectory, which removes between-run
differences. \emph{Alarm level, primary view}: every annotated window counts once; detection is the share
of annotated stagnant windows that drew an alarm within the window plus a two-step tolerance, a false stop
is the same for annotated productive windows, and coverage is the share of \emph{unannotated} windows that
drew one, so silence is visible as zero detection rather than as a low false-stop rate. \emph{Alarm level,
secondary view}: only the first alarm per run counts, scored against that run's annotated regions, with
signed steps saved (a false stop costs the steps it destroys) and the median detection latency in steps.
\emph{Whole-sample}: the alarm rate when the monitor runs on all \pnTraj{} sampled trajectories regardless
of annotation, which shows what an operator would see without labels.
